% !PRINTED 755
\weberpage{755}
%<<CONT
so dass der Grad $n$ von $f(x)$ gleich $(m+1)p-1$ ist, und wir beweisen nun zweierlei:

1) $C_0 F(0) + C_1 F(1) + \cdots + C_m F(m)$ ist eine von Null verschiedene ganze Zahl, also, vom Zeichen abgesehen, mindestens gleich 1,

2) $C_0 Q(0) + C_1 e\,Q(1) + \cdots + C_m e^m Q(m)$ ist kleiner als 1,

beides unter der Voraussetzung, dass über $p$ passend verfügt wird. Ist dies beides bewiesen, so erkennt man die Unmöglichkeit der Gleichung (14) und also die der Gleichung (13), aus der (14) gefolgert war.

Ordnen wir den Zähler von $f(x)$ nach Potenzen von $x$, so ergiebt sich ein Ausdruck der Form:

\weq{16}{f(x) = \frac{A_{p-1}x^{p-1}+A_p x^p + A_{p+1}x^{p+1}+\cdots}{\Pi(p-1)},}

worin $A_{p-1}, A_p, A_{p+1}, \ldots$ ganze Zahlen sind, und $A_{p-1}=[\Pi(m)]^p$, also gewiss nicht durch $p$ theilbar. Es ist daher, wenn man (16) mit der \gesperrt{Taylor}'schen Entwickelung
\[
f(x) = f(0) + xf'(0) + \frac{x^2}{\Pi(2)}f''(0)+\cdots
\]
vergleicht,
\[
f(0)=f'(0)=f''(0)=\cdots=f^{(p-2)}(0)=0,
\]
\[
f^{(p-1)}(0)=A_{p-1},\quad f^{(p)}(0)=pA_p,\quad f^{(p+1)}(0)=p(p+1)A_{p+1}\cdots,
\]
also
\[
F(0) = A_{p-1}+pA_{p+1}+p(p+1)A_{p+1}+\cdots, % [sic] -- printed with A_{p+1} repeated where the pattern would expect A_p, A_{p+1}
\]
eine \gesperrt{durch $p$ nicht theilbare ganze Zahl}. Ordnen wir aber $f(x)$ nach Potenzen von $x-1$, so folgt:
\[
f(x) = \frac{B_p(x-1)^p+B_{p+1}(x-1)^{p+1}+\cdots}{\Pi(p-1)},
\]
worin die $B_p, B_{p+1}, \ldots$ wieder ganze Zahlen sind. Daraus folgt, wie oben, durch Vergleichung mit
\[
f(x) = f(1) + (x-1)f'(1) + \frac{(x-1)^2}{\Pi(2)}f''(1)+\cdots
\]
\[
F(1) = pB_p + p(p+1)B_{p+1}+\cdots,
\]
folglich ist $F(1)$ eine \gesperrt{durch $p$ theilbare ganze Zahl}, und genau auf demselben Wege ergiebt sich, dass $F(2), F(3), \ldots, F(m)$ durch $p$ theilbare ganze Zahlen sind. Da man nun auch $p$ so gross wählen kann, dass $C_0$ nicht durch $p$ theilbar ist, so folgt
%CONT>>
% signature "48*" at foot -- not text, omitted
